\section{Conclusion}
\label{sec:concl}

Runtime bytecode verification at module load time is standard for typed managed virtual machines: the Java Virtual Machine~\cite{jvm} and the Common Language Runtime~\cite{clr} both run a verifier at class-load time to check well-formedness, typing, and stack discipline before bytecode is executed, and the sibling Move VMs deployed elsewhere --- notably the original Diem MoveVM~\cite{OLD_MOVE_LANG} and Sui's variant --- inherit Move's bytecode verifier.
To the best of our knowledge, none of these platforms maintains a \emph{second} layer of runtime safety checks during execution: once the verifier admits a module, its dynamic semantics is trusted.
At the opposite end of the spectrum, dynamically typed languages such as JavaScript or Python perform precisely the runtime type checks of Section~\ref{sec:type}, but lack any first, static layer.
The Ethereum Virtual Machine~\cite{ethereum,kevm} and Solana's SVM sit in a third regime: their bytecode is untyped (a stack of $256$-bit words for the EVM; raw memory regions for the SVM), so neither static bytecode verification nor a runtime type discipline applies.
Aptos's combination of a static bytecode verifier with a complementary runtime layer is, to our knowledge, unique among production VMs and goes beyond the safety guarantees of any other smart-contract platform.

Regarding dynamic reference safety checks, to the best of our knowledge no one has attempted this before.
Existing work on Rust-style borrow semantics~\cite{rustbelt,stackedborrows,treeborrows} is described as a static analysis. The dynamic semantics is also interesting as a reference model for what reference safety means, as most static analyses are over-approximations determined by the particular method chosen, and a canonical semantics does not exist. More work is needed to fully formalize and categorize dynamic reference safety.

Runtime safety checks were added to the Move VM after its original design, which leads to many unnecessary inefficiencies. We are currently working on a new VM that treats dynamic checking and high performance as first principles.
